% !TEX root = ../sop_master/Rhea_SOP.tex
%=======================================
% SOP body, shared by Rhea_SOP.tex and most schools/<school>/ versions.
% Kept short enough that body + Fit stays under 1,000 words.
% Condensed copies (Rhea_SOP_body_800.tex, Rhea_SOP_body_500.tex, and the
% Stanford inline body) are separate -- update them by hand.
%=======================================
\shead{Research interests} I am applying to pursue a PhD in \program{} at \school, where I intend to study questions at the intersection of public and labour economics. I am specifically focused on the effects of public assistance and social program design on labour market outcomes.

\shead{Academic background} I completed a Bachelor of Commerce (Honours) majoring in Economics and Finance at the University of Melbourne, achieving First Class Honours and receiving numerous awards for my academic performance. My Honours thesis, which evaluated the effects of unconventional monetary policy on financial stability in an emerging market context, was also published as a central bank working paper. I have also achieved excellent grades in mathematics subjects, including Real Analysis and Linear Algebra.

\shead{Pre-doctoral and professional experience} After my undergraduate degree, I worked as a senior analyst for three years at the Reserve Bank of Australia, Australia's equivalent of the Federal Reserve. I now work as a predoctoral researcher for Associate Professor Manasi Deshpande on topics in labour and public economics. Our current project examines the labour supply effects of the NDIS, a first-of-its-kind disability support program where the federal government provides individuals with disabilities a budget to spend in the private market on services like personal care, household tasks and physical/behavioural therapy. We exploit the staggered rollout of the NDIS to estimate the causal effects on wages and labour supply. We find that the marginal entrants induced by the NDIS into disability work are low-income, lower-skilled women who obtain modest earnings increases from entry. Our results also suggest that labour responses are heterogeneous by skill level, and skilled occupations such as therapists experience higher wages and lower entry. Through this research I have gained significant experience working with large administrative datasets and independently applying an array of causal inference methods such as staggered rollout event studies, difference-in-differences and shift-share instruments, carefully implemented to ensure credible identification. I believe this has strongly prepared me to conduct independent research at the frontier of the discipline.
\shead{Current research} I am interested in understanding the relative merits of different policy approaches to increase labour supply in high-demand sectors. A common mechanism is to raise wages (e.g., \cite{dalbo_finan_rossi2013}), but its effect on employment is ambiguous in equilibrium, since it would also raise firm labour costs. In addition, if there are high fixed costs to reskilling in these areas, the reservation wage required to compensate for this education may be too high. As such, another potential approach could be to subsidise education, thereby reducing fixed costs of entry \citep{schoonover2026}. The effectiveness of either approach depends on the labour supply elasticity of earnings versus entry costs. \\~\\I believe two policy programs in Australia could help answer this question. Firstly, since 2019 the state of Victoria has implemented a policy known as `Free TAFE', which offers free education in specific courses. This policy was introduced to incentivise students to pursue careers in priority occupations, and as a result had broad eligibility, with residency as the main condition. To contrast this with wage increases, a candidate wage change I could exploit is a national 15 per cent increase for aged care workers from 2023 \citep{fwc_agedcare2022}. Like Free TAFE, this increase is targeted to a specific sector. And unlike other minimum wage changes, the Australian Federal Government is  funding this wage increase by effectively subsidising firms (\href{https://www.health.gov.au/ministers/the-hon-mark-butler-mp/media/budget-2023-24-delivering-the-largest-ever-pay-rise-to-aged-care-workers?language=en&utm_source=miragenews&utm_medium=miragenews&utm_campaign=news}{Australian Government Department of Health (2023)}). This means the relative cost of labour and other inputs is less affected, and so firms may have less incentive to change their input mix. \\~\\
I have commenced this research starting with the education subsidy -- i.e., I identify the causal effect of Free TAFE on enrolment and characterise the marginal enrollee. Understanding who self-sorts is useful to compare against the marginal entrant under minimum wage increases. I do this using Australian administrative data and a difference-in-differences comparing Victoria to other states. I find that Free TAFE generated a large and immediate increase in enrolments, particularly among women, working-aged cohorts, employed workers, and for more advanced courses. Transitioning into individual outcomes such as employment and earnings, I construct a plausible `risk set' of likely enrollees based on ex-ante characteristics and find that young women previously on income support have the strongest enrolment first stage. I have also verified this first stage magnitude using a distance-to-college instrument \citep{card1995}. I am now working on estimating employment and wage responses to Free TAFE. I plan to extend this research using a difference-in-differences to compare entry responses to the aged care minimum wage increase, again by identifying a plausible risk-set.
\shead{Future research} My work on Free TAFE sits within a broader agenda on how labour and general welfare outcomes vary by the different types of support public assistance programs can provide, especially ones that extend beyond income assistance. For instance, another input into the labour supply decision over and above wages and fixed costs of working is family planning (e.g., \cite{gallen2026}). I intend to examine a `free IVF' initiative in Victoria, where IVF is now free for women at a range of public fertility clinics, again with no means-based eligibility tests. This policy effectively lowers one of the `fixed costs' of trying to have a child later, rather than replacing lost income from childbirth. This could generate different childbirth and therefore different labour market decisions than parental income assistance.



%\shead{OLD}

%While minimum wage increases are usually borne by employers \citep{harasztosi_lindner2019}, this was unique in that it was funded by the federal government (and therefore should generate no firm-level changes to labour demand). Ordinarily a wage floor confounds the labour-supply response with firms' labour-demand response, since employers facing a higher wage bill may cut hours or jobs \citep{card_krueger1994,cengiz2019}. Because the Commonwealth financed this rise, labour demand should be unaffected, so I can read the resulting change in aged care employment as a movement along the labour-supply curve and recover an elasticity of entry with respect to the wage. Setting this beside the entry response to Free TAFE's reduction in training costs gives a like-for-like comparison of the two levers: the additional labour each draws into shortage occupations per dollar of public expenditure.

%To explore this question, I plan to use two policy interventions using Australian administrative data. The first is  This policy was introduced to incentivise students to pursue careers in priority occupations, and as a result had broad eligibility (being a Victorian resident was the only condition to apply). Since the cost of courses varies, making the course free effectively represents large education subsidies, which could reduce the fixed cost of labour supply in specific occupations/industries. The second set of policy instruments are raises in the minimum wage in the same industries associated with the `Free TAFE` courses. Plausible candidates would be award increases across [TIMEFRAME], which saw the min wage increase by X. Or, the federal minimum wage increase in the disability sector. 



%The second policy would be inc

%In particular, I would like to compare the effects of targeted education

%I have currently been exploiting an Australian education policy called `Free TAFE', where from 2019 the state of Victoria has offered a free education in certain courses in occupations deemed to have labour shortages.\footnote{Examples include disability care, nursing, construction, and early childhood care. From 2023 onwards, this program was received federal funding to all states.}. This policy was introduced to incentivise students to pursue careers in priority occupations, and as a result had broad eligibility (being a Victorian resident was the only condition to apply). The natural alternative to induce labour market entry would be to increase wages in these occupations. Which lever is appropriate turns on why the shortage persists. If pay does not compensate for the demands of these jobs \citep{rosen1986}---as can happen when employers hold monopsony power \citep{staiger2010}---then raising it should induce entry \citep{clotfelter2008,falch2011}; if instead the up-front cost of the required training is the binding constraint \citep{carneiro_heckman2002}, subsidising it should do more \citep{denning2017,deming_dynarski2009}, though the labour-market returns to such training vary across programs \citep{card_kluve_weber2018,katz2022workadvance}. A subsidy nonetheless earns its keep only if it draws in marginal entrants rather than delivering an inframarginal transfer to those who would have trained anyway \citep{mountjoy2022}. Therefore, two questions therefore arise: firstly, who were the types of cohorts who enrolled due to these courses being free, and did this program successfully lead to persistent employment and earnings increases? And secondly, can we compare the effectiveness of lowering education costs to increasing wages in addressing labour market shortages?
%\\~\\
%I have commenced this research by identifying the causal effect of Free TAFE on enrolments, and documenting the characteristics of additional enrollees. Using Australian administrative data, I construct a triple difference-in-differences comparing enrolments in courses eligible for Free TAFE in Victoria to other states and differences between non free-TAFE eligible courses in Victoria vs other states. So far, I have found that Free TAFE induced higher enrolments among women, working-age cohorts and longer-term, more costly courses such as Certificate IV courses and diplomas. To transition into looking at individual outcomes such as employment and earnings, I have constructed a plausible `risk set' of likely enrollees based on ex-ante characteristics. Specifically, I find that young women previously on income support have the strongest enrolment first stage. I have also verified this first stage using a distance-to-college instrument \citep{card1995}. I will be progressing this work by looking at employment and wage outcomes among this risk set. \\~\\ To answer my second question on comparing wage increases to education costs, 

